\section*{4 Natural Reductions}
We require a technical lemma allowing us to use indirect diagonalization based on robustly-often simulations.
\begin{ntlemma}[Author Lemma9]
Let $T:\mathbb N\to\mathbb N$ be a time-constructible function such that $T(\poly(T(\poly(n))))=o(2^n)$. If SAT is $T(\poly(n))$-robustly often in time $T(n)$, then $\mathrm E\not\subseteq\SIZE(\poly)$.
\end{ntlemma}
\begin{proof}
Suppose SAT is $T(\poly(n))$-robustly often in time $T(n)$. By Proposition3, we have that $\Sigma_2\mathrm P$ is robustly often in time $T(\poly(T(\poly(n))))$. Assume for the sake of contradiction that $\mathrm E\subseteq\SIZE(\poly)$. The Karp--Lipton theorem for $\mathrm E$, credited to Meyer~\cite{ref19}, implies that $\mathrm E$ collapses to $\Sigma_2\mathrm P$. Note that if $T$ is time-constructible, $T'=T(\poly(T(\poly(n))))$ is also time-constructible. Thus we have that $\mathrm E$ is robustly often in time $T'$ for a time-constructible $T'=o(2^n)$, contradicting Proposition4.
\end{proof}
We now re-state and prove Theorem2. In the proof below, we use the fact that SAT is paddable--there is a polynomial-time computable padding function $f$ such that for each $x\in\{0,1\}^*$, and for each $m\geq|x|$, $|f(x,1^m)|=m$ and $f(x,1^m)$ is in SAT iff $x$ is in SAT. Of course, this depends on the encoding of SAT, but it suffices for us that there is some encoding for which paddability is true in the form above, and all other standard properties of SAT continue to hold.
\begin{nttheorem}[Author Theorem10]
If MCSP is NP-hard under natural Turing-reductions, then $\mathrm E\not\subseteq\SIZE(\poly)$.
\end{nttheorem}
\begin{proof}
Suppose that MCSP is NP-hard under natural Turing reductions. Let $M$ be an oracle Turing machine implementing a natural reduction from SAT to MCSP. By the definition of natural reduction, this means that there is a function $s:\mathbb N\to\mathbb N$ such that all queries of $M$ on input of size $n$ for SAT have parameter $s(n)$. If there are no queries that $M$ makes on an input of size $n$, we set $s(n)=0$.
Let $s_1(n)=\log(n)^{\log\log(n)}$, and $s_2(n)=\log(n)^{\log\log\log(n)}$. Moreover let $r(n)=n^{\log(n)}$. We define two Turing machines $M_1$ and $M_2$ that attempt to decide SAT using the oracle machine $M$.
Machine $M_1$ operates as follows. On input $x$ of size $n$, $M_1$ computes strings $x_1\ldots x_{r(n)-n}$, where $x_i=f(x,1^{n+i})$. Here $f$ is the padding function for SAT. Let $x_0=x$. $M_1$ iteratively does the following starting with $i=0$. It simulates $M$ on $x_i$. If an oracle query is made with parameter $k$, it checks if $k$ is at most $s_1(n+i)$. Note that since $M$ implements a natural reduction, $k=s(n+i)$. If $k$ is at most $s_1(n+i)$, it solves the corresponding MCSP instance by brute force search in time at most $\poly(n+i)2^{s_1(n+i)^2}$ and continues the simulation. Note that any succeeding queries must also have parameter at most $s_1(n+i)$, therefore the simulation continues to completion. $M_1$ accepts iff $M$ accepts. If no oracle queries are made, then too the simulation continues to completion, and $M_1$ returns the same answer as $M$.
In case the parameter $k$ for a query is greater than $s_1(n+i)$, $M_1$ increments $i$, and repeats the process, as long as $i\leq r(n)-n$. If $i>r(n)-n$, $M_1$ simply runs the simulation of $M$ on $x$ to completion, now answering all queries by brute-force search regardless of the parameter size.
\begin{ntclaim}
Either SAT is decidable in time $\poly(r(n))2^{s_1(r(n))^2}$, or there is an infinite sequence of disjoint intervals $I_j=[n_j,n'_j)$ of input lengths such that $n'_j=r(n_j)$ for each $j$, and moreover for each $j$, if $n\in I_j$, then $s(n)>s_1(n)$.
\end{ntclaim}
We establish the Claim. Suppose that SAT is not decidable in time $\poly(r(n))2^{s_1(r(n))^2}$. We argue that for any $n_0\in\mathbb N$, there exists an input length $n>n_0$ such that $s(n)\geq s_1(n)$, and moreover $s(m)\geq s_1(m)$ for every $n<m<r(n)$. Indeed, if this were not the case, for every $n>n_0$ one of the iterations of $M_1$ would succeed for some $i<r(n)-n$, and this would mean that the entire simulation would complete in time at most $\poly(r(n))2^{s_1(r(n))^2}$. Note that when the simulation completes, it does solve SAT correctly, as $M$ is an oracle Turing reduction correctly reducing SAT to MCSP, and any completing simulation uses correct answers to oracle queries.
Now we pick the intervals $I_j$ inductively as follows. Let $n_1$ be the least integer such that $s(m)>s_1(m)$ for every $n_1\leq m\leq r(n_1)$, and let $n'_1=r(n_1)$. Inductively, suppose we have picked $n_1\ldots n_j$. We pick $n_{j+1}$ to be the least integer such that $n_{j+1}>r(n_j)$, and moreover $s(m)>s_1(m)$ for each $n_{j+1}\leq m\leq r(n_{j+1})$. Such an integer exists by the argument of the previous paragraph. Let $n'_{j+1}=r(n_j)$, and set $I_j=[n_j,n'_j)$. This sequence of intervals is clearly disjoint, and satisfies the property in the Claim.
The first disjunct of the Claim implies that $\mathrm E\not\subseteq\SIZE(\poly)$, using Lemma9 and the fact that $T(n)=\poly(r(n))n^{s_1(r(n))}$ is time-constructible and satisfies $T(\poly(T(\poly(n))))=o(2^n)$, for the given choice of $r$ and $s_1$. Hence, in the rest of our proof, we assume that the second disjunct holds, i.e., that there is an infinite sequence of long intervals within which each input length generates queries with a large parameter.
Now we define machine $M_2$. On input $x$ of length $n$, $M_2$ uses $M$ to try to solve the search version of SAT, i.e., to find a satisfying assignment for $x$ if one exists. $M_2$ uses the standard search-to-decision reduction for SAT based on self-reducibility, simulating $M$ to answer any decision questions as described below. In order to find a satisfying assignment for $x$, this search-to-decision procedure might call the decision procedure for SAT on inputs $x'$ such that $|x'|<|x|$; in such cases, $M_2$ pads up $x'$ to length $|x|$ by using the padding function $f(x',1^{|x|})$ and runs the simulation of $M$ on $f(x',1^{|x|})$. This guarantees that all inputs on which $M$ is run during the simulation have the same length. At the end of the search-to-decision procedure, if a satisfying assignment is successfully found, $M_2$ accepts, else it rejects.
Next we describe how $M_2$ uses $M$ to decide if an input is in SAT. $M_2$ simulates the oracle machine $M$ on the input. If $M$ asks a query with parameter greater than $s_2(n)$, $M_2$ assumes the answer is ``yes'' and continues the simulation. If $M$ asks a query with parameter at most $s_2(n)$, $M_2$ solves the query by brute force search and continues the simulation. At the end of the simulation, $M_2$ accepts iff $M$ accepts.
$M_2$ always halts within time $\poly(n)2^{s_2(n)^2}$. Unlike machine $M_1$, machine $M_2$ is not guaranteed to solve SAT correctly. However, since $M_2$ has been modified to solve the search version of SAT, the only way $M_2$ can make a mistake on $x$ is to answer ``no'' when the true answer is ``yes''. We will argue that whether or not $M_2$ solves SAT correctly, circuit lower bounds follow.
We consider two cases. The first is that there is a $2^{s_2(n)^3}$-robust set $S$ on which $M_2$ solves SAT correctly. In this case, we can apply Lemma9 with $T(n)=\poly(n)2^{s_2(n)^2}$, noting that $T$ is time-constructible and that $T(\poly(n))<2^{s_2(n)^3}$ for large enough $n$. Thus we get that $\mathrm E\not\subseteq\SIZE(\poly)$ in this case.
The remaining case is that there is no $2^{s_2(n)^3}$-robust set $S$ on which $M_2$ solves SAT correctly. In this case, we apply the main idea of~\cite{ref12}, using in addition our assumption that there is an infinite sequence $\{I_j\}$ of disjoint intervals $[n_j,r(n_j))$ of input lengths such that for each $j$ and $n\in[n_j,r(n_j))$, we have that $s(n)>s_1(n)$.
If there is no $2^{s_2(n)^3}$-robust set $S$ on which $M_2$ solves SAT correctly, there must be a constant $k$ such that for each $n$, there is $m<2^{ks_2(n)^3}$ such that $M_2$ fails to solve SAT correctly on some input of length $m$. Now we use the existence of the sequence $\{I_j\}$ of intervals from the Claim. Since $r(n)=n^{\log(n)}$ is significantly larger than $2^{ks_2(n)^3}$ for large enough $n$, it follows that there is an infinite sequence of input lengths $\{m_i\}$ such that the following hold:(i) $M_2$ fails to solve SAT correctly on some input $x$ of length $m_i$ such that $M$ on input $x$ asks a query with parameter at least $s_1(m_i)$ (ii) For all input lengths $m$ such that $m_i\leq m\leq2^{s_2(m_i)^4}$, $s(m)\geq s_1(m)$.
Now we apply the main idea of~\cite{ref12}, by attempting to use the simulation $M_2$ itself to find a small set of inputs for which $M_2$ fails to decide SAT correctly on at least one of these inputs. We define a Turing machine $A$ as follows. On input $1^m$, $A$ formulates the question of whether there is an instance of size $m$ on which $M_2$ fails to solve SAT correctly and asks a query with parameter at least $s_2(m)$ as an instance $y$ of SAT. Since $M_2$ runs in time $\poly(n)2^{s_2(n)^2}$, such an instance $y$ of size at most $2^{s_2(m)^3}$ for large enough $m$ can be computed in time at most $2^{s_2(m)^3}$. $A$ then runs $M_2$ on $y$ to find a satisfying assignment for $y$, as in the proof of the main result of~\cite{ref12}. For all $m=m_i$ for some $i$, either $A$ succeeds, in which case it finds an instance $x$ of length $m$ on which $M_2$ fails, or it finds a set of at most three instances $y_1,y_2,y_3$ such that $|y_j|=|y|$ for each $1\leq j\leq3$, and $M_2$ fails on at least one of these three instances. Note that by item(ii) in the para above, in the first case $x$ must ask some query with parameter at least $s_1(m)$ that is answered wrongly, and in the second case, one of $y_1,y_2,y_3$ must ask some query with parameter at least $s_1(|y|)$ that is answered wrongly. In either case, a wrongly answered query must be the truth table of a hard function. In the first case, we have that the concatenation $Z_i$ of all queries asked by $x$ is the truth table of a function on $O(\log(m))$ input bits that does not have circuits of size $s_1(m)$, and in the second case we have that the concatenation $Z_i$ of all queries asked by $y_1,y_2,y_3$ must be the truth table of a function on $O(\log(|y|))$ input bits that does not have circuits of size $s_1(|y|)$. $A$ outputs $Z_i$.
In either case, $A$ runs in time $\poly(|Z_i|)2^{s_2(|Z_i|)^3}$ and outputs a string $Z_i$ that does not have circuits of size $s_1(|Z_i|)$. Since $s_2(|Z|)=s_1(|Z|)^{o(1)}$, this implies that $\mathrm E\not\subseteq\SIZE(\poly)$.
\end{proof}
